\section{The unit constraint in the pro-algebra reconstruction}

Throughout this section $k$ is a field and $L/k$ is finite. Tensor
categories and tensor functors are symmetric and unital. Completed
vector spaces and tensor products are the continuous duals of
discrete vector spaces and their coalgebra constructions; ideals
used in a completed quotient are closed.

\begin{lemma}\label{lem:idealizer-takeuchi}
Let $s:L\to A$ be a unital inclusion into a $k$-algebra. In
$D=A\otimes_kA$ let $J$ be the right ideal generated by
$s(\lambda)\otimes1-1\otimes s(\lambda)$, $\lambda\in L$.
Then $D/J$ is the tensor product of the two left $L$-modules $A$.
The quotient of the left idealizer of $J$ is the algebra
\[
 A\times_LA=\left\{\sum a_i\otimes_Lb_i:
 \sum a_is(\lambda)\otimes_Lb_i
 =\sum a_i\otimes_Lb_is(\lambda)\text{ for all }\lambda\right\}.
\]
The statement also holds for the completed products.
\end{lemma}
\begin{proof}
The right ideal imposes exactly the relations
$s(\lambda)a\otimes b=a\otimes s(\lambda)b$. An element $d$
idealizes it precisely when its product on the left with each
displayed generator belongs to $J$, which is the stated invariance
condition. The idealizer is a subalgebra and $J$ is two-sided
inside it. Thus factorwise multiplication descends there, without
assigning multiplication to the entire balanced tensor product.
The same argument with closed relations proves the completed
version, since multiplication is continuous.
This is the equal-source-and-target specialization of Takeuchi's
product~\cite[Section~1, Proposition~3.1 and Definition~4.5]{Takeuchi1977}.
\end{proof}

\begin{example}\label{ex:missing-unit-character}
The algebra and coproduct conditions alone do not force the
character constraint on the counit, even in finite dimension over
$\mathbb F_2$. Let $E=\{a,b,z\}$ be the commutative semigroup
with $a^2=a$, $b^2=b$, $ab=z$, and $z$ absorbing. Its semigroup
algebra has unit $e_a+e_b-e_z$ and the unital augmentation
$e_s\mapsto1$. On its dual $A=k^E$, use pointwise multiplication
and $\Delta f(s,t)=f(st)$. This coproduct is unital,
multiplicative, coassociative and cocommutative. Its coalgebra
counit is $e(f)=f(a)+f(b)-f(z)$, which is not multiplicative:
$e(p_a)e(p_b)=1$ but $e(p_ap_b)=0$.
The proposed tensor category has no allowed unit. Its three
one-dimensional simples satisfy $V_s\otimes V_t=V_{st}$, while
no element of $E$ is a semigroup identity.
\end{example}

\begin{theorem}\label{thm:corrected-pro-bialgebroid-reconstruction}
Let $A$ be a pseudocompact $k$-algebra with continuous multiplication,
source map $s:L\to A$, left $L$-linear continuous counit $e:A\to L$,
and continuous, unital, multiplicative, cocommutative and
coassociative coproduct into the completed Takeuchi product.
Assume the scalar and ordinary counit compatibilities, and impose
\[
                   e(1)=1,\qquad e(ab)=e(a\,s(e(b))).
\]
These are the data of a continuous cocommutative left
$L$-bialgebroid with equal source and target. Its finite-dimensional
continuous modules form a symmetric abelian tensor category with
exact faithful fibre functor to finite-dimensional $L$-spaces.
Conversely an essentially small abelian $k$-linear symmetric
tensor category with such a fibre functor reconstructs these data,
and the two constructions are inverse up to equivalence respecting
the fibre functor and the induced continuous pro-algebra isomorphism.
\end{theorem}
\begin{proof}
Define the action on $M\otimes_LN$ by the coproduct. The Takeuchi
condition makes it balanced, and all sums act through finite
quotients. Multiplicativity gives the action law; coassociativity
and cocommutativity give the usual associativity and symmetry.
The unit object is $L$ with action
$a\cdot\lambda=e(a s(\lambda))$. The character equation gives
the module law, while $e(1)=1$ and left $L$-linearity give the
unit law and the prescribed scalar action. Apply the counit
identity to $\Delta(a s(\lambda))$ to check that both ordinary
unit constraints are $A$-linear. Coherence is the ordinary linear
coherence. Kernels and cokernels remain finite and continuous,
and tensoring over the field $L$ is exact.

For the converse let $F$ be the fibre functor, set
$W=\operatorname{Res}_{L/k}F$, and take
$A=\operatorname{Nat}_k(W,W)$ with its finite topology. Hom spaces
are finite-dimensional over $k$, since $L/k$ is finite and $F$ is
faithful. Exact faithfulness bounds the length of a strict chain
of subobjects by $\dim_LF(M)$. Thus the category is locally finite.
The coalgebra reconstruction theorem identifies it with finite
continuous $A$-modules; its hypotheses are precisely those just
checked~\cite[6.1--6.2]{Deligne1990}\cite[Theorem~2.1.2]{Hai2008}.
Scalar multiplication gives $s:L\to A$.

A natural endomorphism of $(M,N)\mapsto M\otimes_LN$ is determined
on pairs of finite regular quotients $A/I,A/J$ by its value at
$1\otimes1$. The maps $a\mapsto am$ then determine its value on
every pair of modules. Naturality for right multiplication by
$s(\lambda)$ is exactly the Takeuchi condition. Compatible values
on all such quotients are consequently the completed Takeuchi
product. The action on tensor-product objects reconstructs
$\Delta$; the action on the unit reconstructs $e$. The action law,
associativity and symmetry give all the asserted equations,
including the character identity. Continuous dualization gives
the commutative unital $L$-algebra dual to the coproduct; the
transpose of $1_A$ is its augmentation. The same regular-quotient
argument recovers $A$ as the natural endomorphism algebra, proving
the inverse comparison. The coend formulation is
Hai's~\cite[Corollary~2.2.5]{Hai2008}; no finite-dimensional dual
is substituted for the continuous dual without a finiteness proof.
\end{proof}

\begin{proposition}\label{prop:restricted-diagonal-fibre-functor}
If the preceding category is rigid and $k=\operatorname{End}(1)$,
its scalar extension to $L$ has the exact faithful fibre functor
induced by $F$. In the pro-algebra description this is the diagonal
functor on the corresponding finite continuous
$(A\widehat\otimes_kL)$-modules. The conclusion includes inseparable
$L/k$. The diagonal functor on unrestricted
$(L\otimes_kL)$-modules need not be exact and need not be faithful.
\end{proposition}
\begin{proof}
Deligne identifies the category with representations of the affine
groupoid over $\operatorname{Spec}L$ whose arrow scheme is faithfully
flat over $\operatorname{Spec}(L\otimes_kL)$
\cite[Theorem~1.12]{Deligne1990}. After scalar extension, restrict
an equivariant vector bundle to the diagonal $L$-point. Local
freeness makes this exact; transitivity makes it faithful, since
after the faithfully flat base change of arrows from that point,
every equivariant map is transported from its value there.
Deligne's scalar-extension compatibility identifies this with
$M\mapsto M\otimes_{L\otimes_kL}L$
\cite[5.18--5.21]{Deligne1990}.

For a purely inseparable extension of degree $p$,
$L\otimes_kL=L[\delta]/(\delta^p)$. Tensoring the injection
$(\delta^{p-1})\hookrightarrow L[\delta]/(\delta^p)$ with the
diagonal quotient $L$ gives a nonzero source and zero map.
For a nontrivial separable extension the diagonal component is an
idempotent factor; the other factor is nonzero and is killed.
These are unrestricted-module counterexamples, not objects with
the required equivariance.
\end{proof}

\begin{proposition}\label{prop:continuous-dual-and-rigidity}
Choose $0\ne\ell:L\to k$ linear, and put
$U=\operatorname{Hom}^{\rm cont}_L(A,L)$, with
$(au)(b)=u(ba)$. Then
\[
 \operatorname{Hom}_A(M,U)\simeq\operatorname{Hom}_L(M,L),\qquad
 \operatorname{End}_A(U)\simeq A^{\rm op}.
\]
A rigid tensor structure induces a continuous anti-isomorphism
fixing $s(L)$ and hence a second commuting action on $U$.
Conversely such an anti-isomorphism gives rigidity exactly when,
for every finite continuous module, the ordinary evaluation and
coevaluation maps are $A$-linear. Their triangle identities then
hold automatically. Testing only the infinite pro-module $A$
does not replace the finite-module tests.
\end{proposition}
\begin{proof}
Evaluation at $1$ is the first map; its inverse sends $f$ to
$\phi_f(m)(b)=f(bm)$. The pairing $\ell(xy)$ on $L$ is
nondegenerate: for $x\ne0$ choose $z$ with $\ell(z)\ne0$ and
take $y=x^{-1}z$. Thus $u\mapsto\ell\circ u$ identifies $U$
with the continuous $k$-dual of $A$, compatibly with the action.
Continuous duality identifies its endomorphism ring with the
opposite of the continuous right-regular endomorphism ring of $A$,
which is $A^{\rm op}$ by evaluation at $1$.

For a categorical dual $DM$, identify $F(DM)$ with the $L$-dual
of $F(M)$ and use the nondegenerate $k$-pairing
$\ell(\phi(m))$. Take the adjoint of the action on $DM$.
Naturality gives an element of $A$, transposition reverses
products, and the bidual isomorphism gives the inverse. Scalars
are self-adjoint, so the map fixes $s(L)$; the construction is
continuous in the finite topology. It realizes the dual module
through $\operatorname{Hom}_A(M,U)$.

Conversely use the anti-isomorphism and the commuting action to
give that Hom space its dual action. On underlying spaces this
is ordinary exact $L$-duality. With dual bases $m_i,\phi_i$,
the two maps are evaluation and $1\mapsto\sum_i m_i\otimes\phi_i$.
The triangle identities are
$\sum_i m_i\phi_i(m)=m$ and
$\sum_i\phi(m_i)\phi_i=\phi$. They give categorical duality
precisely when both maps are $A$-linear. Explicitly these conditions
are
\[
 \sum\operatorname{ev}(a_{(1)}m\otimes a_{(2)}\phi)
       =e(a s(\operatorname{ev}(m,\phi))),
 \quad
 \sum_{i,(a)}a_{(1)}m_i\otimes a_{(2)}\phi_i
       =e(a)\sum_i m_i\otimes\phi_i.
\]
For separable $L/k$ the trace is a permissible canonical choice of
$\ell$. In a purely inseparable extension the trace is zero; the
argument uses an arbitrary nonzero functional instead.
\end{proof}

\paragraph{Comparison with established structures.}\label{par:qualified-tannaka-comparison}
Takeuchi's Definition~4.5 requires both the coproduct and the map
to $\operatorname{End}_k(L)$ to be algebra maps. The latter map is
$a\mapsto(\lambda\mapsto e(a s(\lambda)))$, so its multiplicativity
is exactly the missing character equation. Schauenburg's
Theorem~5.1 and Remark~5.3 identify bialgebroid structures with
monoidal module structures and permit transport from a non-strict
underlying functor~\cite{Schauenburg1998}; the regular module in
that proof becomes the family of finite regular quotients here.
Hai's coend theorems provide the abelian reconstruction and
rigid refinements~\cite[Theorems~2.1.2, 2.2.4, 2.2.6 and
Corollary~2.2.5]{Hai2008}. The weak-bialgebra reformulation requires
a Frobenius-separable base~\cite[Section~3 and Theorems~5.1, 5.5]{Schauenburg2003};
a finite separable field extension has that property, while a
purely inseparable extension does not. This restriction is distinct
from the unrestricted reconstruction and restricted diagonal
statements. These comparisons make no priority or independence claim.

\section{Analysers and the extra plethory data}

\begin{proposition}\label{prop:divided-power-analyser}
The positive divided-power lattice
$\Gamma_+=\bigoplus_{n\ge1}\mathbb Z\gamma_n\subset\mathbb Q[T]$,
$\gamma_n=T^n/n!$, is closed under multiplication and composition,
with composition identity $T$. It acts on every divided-power
ideal, including ideals in rings with torsion. Its constant ring
is zero. Its additive decomposition does not take values in
$\Gamma_+\otimes\Gamma_+$; it takes values in the augmentation
ideal of the tensor product of its unitizations.
\end{proposition}
\begin{proof}
Multiplication has coefficients $\binom{m+n}m$. If
$G=\sum a_i\gamma_i$, the coefficient of $\gamma_N$ in
$\gamma_m(G)$ indexed by $\sum r_i=m$, $\sum ir_i=N$ is
\[
                 \frac{N!}{\prod_i r_i!(i!)^{r_i}}
                              \prod_i a_i^{r_i}.
\]
The integer factor counts set partitions with the prescribed block
sizes. These identities follow from the divided-power axioms
without division in the coefficient ring, so they give the
asserted action even with torsion. The ordinary polynomial
identities in $\mathbb Q[T]$ give composition associativity.
The additive formula is
\[
 \gamma_n(X+Y)=\gamma_n(X)+\gamma_n(Y)
                    +\sum_{i=1}^{n-1}\gamma_i(X)\gamma_{n-i}(Y),
\]
while $\gamma_n(XY)=n!\gamma_n(X)\gamma_n(Y)$.
Every factor in $\Gamma_+$ vanishes at zero, so a sum of products
from $\Gamma_+\otimes\Gamma_+$ would vanish when $Y=0$; the
left side does not. This proves the additive qualification and
also the failure of faithful evaluation on the zero constant ring.
\end{proof}

\begin{theorem}\label{thm:analyser-to-plethory}
For a unital analyser $P$ with specified constant subring $k$, a
compatible $k$-plethory consists exactly of a $k$--$k$ biring
structure realizing its additive and multiplicative decompositions,
with scalar evaluation $\beta(c)(F)=F\circ c$, such that composition
$P\odot_kP\to P$ and the unit $k[T]\to P$ are biring morphisms.
Existence of unspecified finite decompositions does not impose the
biring coherence identities. If the full constant ring is a field
$k$, evaluation on it is faithful, and the unital analyser axioms
and both decomposition axioms hold, there is a unique compatible
$k$-plethory.
\end{theorem}
\begin{proof}
The first assertion unpacks the monoid definition of a plethory
and the relations of the composition product
\cite[1.1, 1.3, 1.8]{BorgerWieland2005}. Linearity and
multiplicativity in the outer argument, the two decomposition
relations in the inner argument, and the scalar relation are
exactly what allows composition to descend. The remaining
condition is that composition and the unit respect the biring maps.

For the sufficient criterion regard $P$ as a subalgebra of $k^k$.
Evaluation embeds $P^{\otimes_k n}$ into $k^{k^n}$ for every $n$.
For two factors, a finite independent family of functions in the
first factor has evaluation vectors spanning its dual; sufficiently
many evaluations recover each coefficient function. Induction
gives all $n$. The two decompositions therefore define unique
coproducts. Every biring identity can be checked on constants,
where it is an ordinary field identity, and injectivity makes
the identity hold in the tensor product. The counits are
evaluation at $0,1$, the coinverse is composition with $-T$,
and the scalar map is evaluation at constants. Composition
respects each structure by the same evaluation test; its
associativity and unit are the analyser identities.
\end{proof}

\begin{proposition}\label{prop:constant-evaluation-criterion}
For a polynomial analyser $\Omega(R,S)$ the constants are exactly
$R$, and faithful evaluation on them is equivalent to
\[
 \bigcap_{r\in R}\ker\bigl(S^{d+1}\to S,
                   (a_i)\mapsto\sum_i a_ir^i\bigr)=0
                            \quad(d\ge0).
\]
If $S$ is a domain this holds exactly when $R$ is infinite.
More generally it holds whenever, for each $d$, there is a
$(d+1)$-point Vandermonde determinant regular in $S$.
\end{proposition}
\begin{proof}
A constant under composition equals its composition with zero,
so is an ordinary constant in $R$. The kernel of evaluation is
the null-polynomial ideal, giving the coefficient criterion.
A nonzero polynomial over a domain has at most its degree many
roots, while for finite $R$ the monic product
$\prod_{r\in R}(T-r)$ is a nonzero null polynomial. The adjugate
of a Vandermonde matrix annihilates each null coefficient by its
determinant, proving the regular-determinant implication.
\end{proof}

\begin{example}\label{ex:finite-analyser-without-biring}
The ring $P=\mathbb F_2[T]/((T^2-T)^3)$ has $64$ elements and
is a unital analyser satisfying both finite decomposition axioms,
but admits no compatible plethory structure.
\end{example}
\begin{proof}
For every polynomial $G$, $G^2-G$ vanishes at $0,1$ and is
divisible by $T^2-T$. Consequently substitution preserves the
outer defining relation. Changing the inner representative also
preserves it, by the polynomial difference identity. Ordinary
substitution therefore descends, with its unit, associativity and
outer ring laws. Binomial and diagonal monomial expansions give
the two finite decomposition formulas. The constants are
$\mathbb F_2$, by composing with zero.
Put $q=T^2-T$. A compatible unit forces $q$ to be primitive for
coaddition, so
\[
                  \Delta_+(q^3)=q^2\otimes q+q\otimes q^2.
\]
This is nonzero because $q,q^2$ are linearly independent, whereas
$q^3=0$ in $P$. Hence a compatible coaddition cannot be an algebra
map. This also explains why the faithful-evaluation hypothesis
in the preceding theorem is necessary.
\end{proof}
